% =============================================================================
% Chapter 10: APB Peripheral Testbench Guide
% =============================================================================
\chapter{APB Peripheral Testbench Guide}
\label{ch:apb_guide}

This chapter extends the SPI master verification methodology to other APB peripherals in the same
design family, including APB Timer, APB GPIO, APB UART, and similar blocks. The goal is to provide
a reusable framework that minimizes testbench development time for each new peripheral.

\section{APB Protocol Fundamentals}

All peripherals in this SoC use the AMBA 3 APB (Advanced Peripheral Bus) protocol. Understanding
the protocol is essential for building correct testbenches.

\subsection{APB Signal Definitions}

\begin{table}[H]
\centering
\caption{APB Signal Definitions (Common to All APB Peripherals)}
\label{tab:apb_signals_def}
\begin{tabularx}{\textwidth}{@{}lclX@{}}
\toprule
\textbf{Signal} & \textbf{Direction} & \textbf{Width} & \textbf{Description} \\
\midrule
\texttt{PCLK / HCLK}  & Input  & 1  & APB clock (system clock). All transfers are synchronous to this. \\
\texttt{PRESETn / HRESETn} & Input & 1 & Active-low reset. All registers clear on assertion. \\
\texttt{PADDR}   & Input  & N  & Byte address of the target register. Width varies by peripheral (typically 8--12 bits). \\
\texttt{PSEL}    & Input  & 1  & Peripheral select. Asserted for 2 clock cycles per transfer. \\
\texttt{PENABLE} & Input  & 1  & Access phase enable. Low during setup phase, high during access phase. \\
\texttt{PWRITE}  & Input  & 1  & Transfer direction: 1 = write, 0 = read. \\
\texttt{PWDATA}  & Input  & 32 & Write data from master to slave. \\
\texttt{PRDATA}  & Output & 32 & Read data from slave to master. \\
\texttt{PREADY}  & Output & 1  & Slave ready. 1 = transfer completes this cycle. \\
\texttt{PSLVERR} & Output & 1  & Slave error. 1 = transfer error (rarely used). \\
\bottomrule
\end{tabularx}
\end{table}

\subsection{APB Transfer Timing}

Every APB transfer takes exactly two phases:

\begin{enumerate}[leftmargin=2em]
  \item \textbf{Setup Phase} (1 clock cycle): Master asserts \texttt{PSEL=1}, \texttt{PENABLE=0},
        drives \texttt{PADDR}, \texttt{PWRITE}, and \texttt{PWDATA} (for writes).
  \item \textbf{Access Phase} (1 clock cycle): Master asserts \texttt{PENABLE=1}. The slave
        samples write data or provides read data. If \texttt{PREADY=1}, the transfer completes.
        If \texttt{PREADY=0}, the access phase extends until \texttt{PREADY} goes high.
\end{enumerate}

\begin{verbatim}
PCLK:    _/‾\_/‾\_/‾\_/‾\
PSEL:    ____/‾‾‾‾‾‾‾\____
PENABLE: ________/‾‾‾\____
PADDR:   ----<ADDR>--------
PWDATA:  ----<DATA>--------   (write)
PRDATA:  --------<DATA>----   (read, valid when PENABLE & PREADY)
PREADY:  ‾‾‾‾‾‾‾‾‾‾‾‾‾‾‾‾   (always 1 for simple peripherals)
\end{verbatim}

\subsection{APB Read Timing Considerations}

For peripherals with FIFO read ports (SPI, UART, etc.), the read data must be sampled at the
clock edge during the access phase, \emph{before} any simulation delay (\texttt{\#1}). This is
because the FIFO read pointer advances on the same edge via non-blocking assignment.

For peripherals with simple register reads (Timer, GPIO), this timing is less critical because
the register value does not change as a side effect of reading.

\section{Reusable Testbench Template}

The following template shows the common structure for any APB peripheral testbench. Only the
peripheral-specific sections (marked with comments) need to change between peripherals.

\begin{lstlisting}[caption={Reusable APB Peripheral Testbench Template},label={lst:template}]
`timescale 1ns/1ps
module tb_apb_<peripheral>;

    parameter APB_ADDR_WIDTH = 12;
    parameter CLK_PERIOD = 10;

    // === APB signals (COMMON - same for all peripherals) ===
    logic                      HCLK, HRESETn;
    logic [APB_ADDR_WIDTH-1:0] PADDR;
    logic               [31:0] PWDATA, PRDATA;
    logic                      PWRITE, PSEL, PENABLE;
    logic                      PREADY, PSLVERR;

    // === Peripheral-specific signals (CHANGE PER IP) ===
    // For Timer: logic [31:0] timer_val; logic timer_irq;
    // For GPIO:  logic [31:0] gpio_in, gpio_out, gpio_dir;
    // For UART:  logic uart_tx, uart_rx;
    // For SPI:   logic spi_clk, spi_csn*, spi_sdo*, spi_sdi*;

    // === DUT instantiation (CHANGE PER IP) ===
    apb_<peripheral> #(...) dut (...);

    // === Clock generation (COMMON) ===
    initial HCLK = 0;
    always #(CLK_PERIOD/2) HCLK = ~HCLK;

    // === Pass/fail counters (COMMON) ===
    int pass_cnt = 0, fail_cnt = 0;

    // === APB BFM tasks (COMMON - reuse without changes) ===
    function automatic logic [APB_ADDR_WIDTH-1:0]
        reg_addr(input logic [3:0] idx);
        return {6'h0, idx, 2'b00};
    endfunction

    task automatic apb_write(...);  // Same as SPI testbench
    task automatic apb_read(...);   // Same as SPI testbench
    task automatic check(...);      // Same as SPI testbench

    // === Peripheral-specific BFM (CHANGE PER IP) ===
    // For Timer: no BFM needed, just wait for interrupts
    // For GPIO:  drive gpio_in, monitor gpio_out
    // For UART:  serial TX/RX BFM tasks
    // For SPI:   slave BFM (see Chapter 3)

    // === Test sequence (CHANGE PER IP) ===
    initial begin
        // Reset
        PADDR = '0; PWDATA = '0; PWRITE = 0; PSEL = 0; PENABLE = 0;
        HRESETn = 0;
        repeat(5) @(posedge HCLK);
        HRESETn = 1;
        repeat(3) @(posedge HCLK);

        // TC1: Reset check (COMMON pattern)
        // TC2: Register read-back (COMMON pattern)
        // TC3+: Functional tests (PERIPHERAL SPECIFIC)

        // Summary (COMMON)
        $display("PASSED: %0d / FAILED: %0d", pass_cnt, fail_cnt);
        $finish;
    end

    // === Watchdog (COMMON) ===
    initial begin #10_000_000; $display("TIMEOUT!"); $finish; end

    // === Waveform dump (COMMON) ===
    initial begin $recordfile("waves.trn"); $recordvars(); end
endmodule
\end{lstlisting}

\section{Applying the Template to APB Timer}

An APB Timer peripheral typically has:
\begin{itemize}[leftmargin=2em]
  \item A control register (enable, mode, prescaler).
  \item A compare/match register.
  \item A counter register (read-only or read/write).
  \item An interrupt status/clear register.
  \item An interrupt output signal.
\end{itemize}

\subsection{Timer Testcase Structure}

\begin{table}[H]
\centering
\caption{Example Testcases for APB Timer}
\label{tab:timer_tc}
\begin{tabularx}{\textwidth}{@{}clX@{}}
\toprule
\textbf{TC} & \textbf{Name} & \textbf{What to Verify} \\
\midrule
TC1 & Reset & Counter = 0, control = 0, interrupt = 0 after reset. \\
TC2 & Register R/W & Write prescaler, compare value; read back and verify. \\
TC3 & Basic count & Enable timer, wait N cycles, read counter, verify it incremented. \\
TC4 & Compare match & Set compare value, enable timer, wait for interrupt, verify timing. \\
TC5 & Prescaler & Set prescaler, verify counter increments at divided rate. \\
TC6 & Interrupt clear & After match interrupt, write clear register, verify interrupt deasserts. \\
TC7 & Reload/wrap & Verify counter wraps or reloads after reaching compare value. \\
TC8 & Disable & Disable timer mid-count, verify counter stops. \\
\bottomrule
\end{tabularx}
\end{table}

\subsection{Timer Signal Observation}

Monitor the same APB bus signals as SPI, plus:
\begin{itemize}[leftmargin=2em]
  \item Timer counter value (hierarchical access to internal counter register).
  \item Compare match flag.
  \item Interrupt output.
  \item Prescaler counter (to verify divided clock).
\end{itemize}

\section{Applying the Template to APB GPIO}

An APB GPIO peripheral typically has:
\begin{itemize}[leftmargin=2em]
  \item Direction register (input/output per pin).
  \item Output data register.
  \item Input data register (read-only).
  \item Interrupt enable, edge-select, and status registers.
\end{itemize}

\subsection{GPIO Testcase Structure}

\begin{table}[H]
\centering
\caption{Example Testcases for APB GPIO}
\label{tab:gpio_tc}
\begin{tabularx}{\textwidth}{@{}clX@{}}
\toprule
\textbf{TC} & \textbf{Name} & \textbf{What to Verify} \\
\midrule
TC1 & Reset & All outputs tri-stated or low, direction = input. \\
TC2 & Register R/W & Write direction and output registers; read back. \\
TC3 & Output drive & Set pin as output, write value, verify on gpio\_out port. \\
TC4 & Input read & Drive gpio\_in externally, read input register, verify match. \\
TC5 & Direction switch & Change pin from input to output and back; verify behavior. \\
TC6 & Edge interrupt & Enable rising-edge interrupt, toggle gpio\_in, verify interrupt fires. \\
TC7 & Interrupt clear & After interrupt, write clear, verify interrupt deasserts. \\
TC8 & Multiple pins & Verify independent operation of multiple GPIO pins simultaneously. \\
\bottomrule
\end{tabularx}
\end{table}

\subsection{GPIO BFM}

The GPIO BFM is simpler than the SPI slave BFM---it just drives the \texttt{gpio\_in} bus:
\begin{lstlisting}
logic [31:0] gpio_in_drive = '0;
assign gpio_in = gpio_in_drive;
// In testcase: gpio_in_drive = 32'h0000_00A5;
\end{lstlisting}

\section{Applying the Template to APB UART}

An APB UART peripheral typically has:
\begin{itemize}[leftmargin=2em]
  \item TX data register (write to transmit).
  \item RX data register (read to receive).
  \item Baud rate divisor register.
  \item Status register (TX empty, RX full, overrun, etc.).
  \item Control register (enable, parity, stop bits).
  \item TX/RX FIFOs.
\end{itemize}

\subsection{UART Testcase Structure}

\begin{table}[H]
\centering
\caption{Example Testcases for APB UART}
\label{tab:uart_tc}
\begin{tabularx}{\textwidth}{@{}clX@{}}
\toprule
\textbf{TC} & \textbf{Name} & \textbf{What to Verify} \\
\midrule
TC1 & Reset & TX idle (high), status shows empty, no errors. \\
TC2 & Register R/W & Write baud divisor, control register; read back. \\
TC3 & TX single byte & Write byte to TXFIFO, verify serial output matches. \\
TC4 & RX single byte & Drive serial input with known byte, read RXFIFO, verify. \\
TC5 & Loopback & Connect TX to RX, send byte, verify receive. \\
TC6 & Baud rate & Verify bit timing matches baud rate divisor setting. \\
TC7 & FIFO depth & Fill TX FIFO, verify backpressure behavior. \\
TC8 & RX interrupt & Enable RX interrupt, receive byte, verify interrupt fires. \\
\bottomrule
\end{tabularx}
\end{table}

\subsection{UART BFM}

A UART serial BFM is needed to drive \texttt{uart\_rx} with properly timed serial data:
\begin{lstlisting}
task uart_send_byte(input logic [7:0] data);
    integer i;
    uart_rx = 1'b0;  // Start bit
    #(baud_period);
    for (i = 0; i < 8; i++) begin
        uart_rx = data[i];  // LSB first
        #(baud_period);
    end
    uart_rx = 1'b1;  // Stop bit
    #(baud_period);
endtask
\end{lstlisting}

\section{Common Testcase Patterns Across All APB Peripherals}

The following testcase patterns are universal and should be included in every APB peripheral
testbench:

\subsection{Pattern 1: Reset Verification}

After releasing reset, verify:
\begin{itemize}[leftmargin=2em]
  \item All control registers are at their default values.
  \item All status registers show idle/safe state.
  \item All output ports are in their default (safe) state.
  \item No interrupts are asserted.
  \item \texttt{PSLVERR} is 0.
\end{itemize}

\subsection{Pattern 2: Register Read-Back}

For every read/write register:
\begin{itemize}[leftmargin=2em]
  \item Write a known pattern (e.g., 0xA5A5A5A5, 0x5A5A5A5A, 0xFFFFFFFF, 0x00000000).
  \item Read back and verify each bit field.
  \item Check that read-only fields are not affected by writes.
  \item Check that reserved fields read as zero.
\end{itemize}

\subsection{Pattern 3: Basic Functional Operation}

Perform the simplest possible functional operation:
\begin{itemize}[leftmargin=2em]
  \item Timer: Enable, wait, check counter > 0.
  \item GPIO: Set output, check pin.
  \item UART: Send byte, check serial output.
  \item SPI: Write data, check EOT.
\end{itemize}

\subsection{Pattern 4: Interrupt Verification}

For every interrupt source:
\begin{enumerate}[leftmargin=2em]
  \item Verify no interrupt before enabling.
  \item Enable the interrupt.
  \item Create the trigger condition.
  \item Verify interrupt asserts.
  \item Clear the interrupt.
  \item Verify interrupt deasserts.
  \item Verify no spurious re-assertion.
\end{enumerate}

\subsection{Pattern 5: Software Reset}

If the peripheral supports software reset:
\begin{itemize}[leftmargin=2em]
  \item Configure the peripheral.
  \item Load FIFOs or start an operation.
  \item Issue software reset.
  \item Verify the peripheral returns to the post-reset state.
\end{itemize}

\section{Signal Observation Method (Common to All APB Peripherals)}

The same signal-observation methodology used for the SPI master applies to all APB peripherals:

\begin{enumerate}[leftmargin=2em]
  \item \textbf{APB bus signals:} PADDR, PWDATA, PRDATA, PWRITE, PSEL, PENABLE---always monitor
        these to verify correct register addressing and data transfer.
  \item \textbf{Control registers:} Use hierarchical access to monitor internal register values
        and verify they match APB writes.
  \item \textbf{Status flags:} Monitor status register bits to track peripheral state transitions.
  \item \textbf{Output ports:} Verify that physical outputs (SPI clock, GPIO pins, UART TX)
        match expected behavior.
  \item \textbf{Interrupt outputs:} Monitor for correct assertion, deassertion, and absence of
        spurious pulses.
  \item \textbf{Internal counters:} Use analog waveform display for counters to quickly verify
        counting behavior.
  \item \textbf{FIFO levels:} Monitor element counts to verify push/pop operations.
\end{enumerate}

\section{Reusability Checklist}

When creating a testbench for a new APB peripheral, follow this checklist:

\begin{enumerate}[leftmargin=2em]
  \item[$\square$] Copy the template from Section~10.2.
  \item[$\square$] Update the DUT instantiation and peripheral-specific signals.
  \item[$\square$] Update the register address defines.
  \item[$\square$] Implement TC1 (reset check) using the common pattern.
  \item[$\square$] Implement TC2 (register read-back) for all registers.
  \item[$\square$] Implement functional testcases specific to the peripheral.
  \item[$\square$] Implement interrupt testcases using the common pattern.
  \item[$\square$] Add hierarchical probes for key internal signals.
  \item[$\square$] Run with waveform dump enabled.
  \item[$\square$] Verify all checks pass.
  \item[$\square$] Document results in a report following this template.
\end{enumerate}
